\subsection{Mixed-characteristic specialization and the branch character at a node}

The specialization question on 67/049(c) does not exclude a positive
residue characteristic. We therefore distinguish the full-cover statement
over a characteristic-zero trait from the following tame statement.
For a log-regular model $\mathcal X^{\log}$ with interior $U$, write
$\operatorname{Tame}_{\mathcal X}(U)$ for the full category of finite
\'etale covers of $U$ whose normalization over $\mathcal X$ is tamely
ramified at every codimension-one point of
$\mathcal X\setminus U$. Tameness here includes separability of the
residue extension as well as ramification index prime to the residue
characteristic; we use exactly the condition in
\cite[Proposition~A.10, pp.~57--58]{HoshiLog2009A567}.

\begin{theorem}[Known comparison with proved model independence:
 mixed-characteristic tame specialization]\label{A6-tame-specialization}
Let $R$ be a complete noetherian discrete valuation ring, with fraction
field $F$ of characteristic zero and residue field $k$ of characteristic
$p\geq0$. Give $S=\operatorname{Spec}R$ the divisorial log structure of
its closed point. Let $U/F$ be a smooth geometrically connected
separated finite-type variety.
Suppose that $\mathcal X^{\log}\to S^{\log}$ is a proper morphism of
noetherian fs log schemes, that $\mathcal X^{\log}$ is integral and log
regular, and that its interior is identified with $U$. Write
$\mathcal X_0^{\log}$ for its strict closed log fibre. Then restriction
gives equivalences
\[
 \operatorname{Tame}_{\mathcal X}(U)
 \simeq\operatorname{K\acute Et}(\mathcal X^{\log})
 \simeq\operatorname{K\acute Et}(\mathcal X_0^{\log}).
\]
They commute with pullback from $S^{\log}$.

If $\mathcal Y^{\log}$ is any other model satisfying these same
hypotheses for the same $U$ and $R$, then
\[
 \operatorname{Tame}_{\mathcal X}(U)
       =\operatorname{Tame}_{\mathcal Y}(U)
       \quad\text{as full subcategories of }\operatorname{F\acute Et}(U).
\]
Consequently the special-log-fibre cover categories are canonically
equivalent through this common category. The comparisons are independent
of a common refinement, whose existence is not assumed. They are unique
up to the natural isomorphism normalized by their restriction to $U$.
For an actual log morphism of models fixing $U$, its pullback is this
equivalence.

Let $G_F^t$ be the quotient of $G_F$ classifying finite extensions tame
over $R$, and let $\Pi^t(U)$ be the fundamental group of the common
tame-cover category. There is a natural surjection
$\Pi^t(U)\to G_F^t$. A uniformizer $\varpi$ supplies a section
\[
 \sigma_{\varpi}:G_k\longrightarrow G_F^t,
 \qquad
 G_F^t\simeq\widehat{\Z}^{(p')}(1)\rtimes G_k,
\]
where all primes are retained when $p=0$. If
$\Delta^t=\ker(\Pi^t(U)\to G_F^t)$, the resulting tangential extension
\[
 1\longrightarrow\Delta^t\longrightarrow
 \Pi^t(U)\times_{G_F^t}G_k
 \longrightarrow G_k\longrightarrow1
\]
is independent of the chosen model. For $p>0$ we do not identify
$\Delta^t$ with the full geometric \'etale fundamental group, nor do
we replace $G_F^t$ by $G_F$.
\end{theorem}
\begin{proof}
The first equivalence is the stated Proposition~A.10. The second is
\cite[Corollary~3.9, pp.~31--32]{HoshiLog2009A567}, applied to the proper
log-regular total space over the complete local base. Neither source
requires the residue field to have characteristic zero. Their functors
are the restrictions to the interior and to the strict closed log
fibre; hence their compatibility with base pullback is literal
functoriality. More generally, the same two equivalences hold over any
complete noetherian local fs log base for an existing proper log-regular
total space, with its own interior and the indicated boundary-tame
category. The trait hypotheses are used below for the arithmetic
description and its frame.

We prove model independence without resolution. A log-regular scheme
is normal, and its log structure is its divisorial log structure from
the interior; this is the fact used in the proof of Proposition~A.10
on p.~58. Take $V\to U$ in
$\operatorname{Tame}_{\mathcal X}(U)$ and extend it to its finite
Kummer-\'etale cover of $\mathcal X^{\log}$. Let $\eta$ be a
codimension-one point of $\mathcal Y\setminus U$ and put
$A=\mathcal O_{\mathcal Y,\eta}$. Normality makes $A$ a discrete
valuation ring, with fraction field the common function field $F(U)$.
The generic-point map
$\operatorname{Spec}F(U)\to U\hookrightarrow\mathcal X$ extends to
$\operatorname{Spec}A\to\mathcal X$ by properness over $S$.
Give $\operatorname{Spec}A$ the divisorial log structure of its closed
point. The extension is a log morphism: a local boundary-log section
on $\mathcal X$ is a regular function invertible on $U$, whose pullback
is a nonzero element of $A$. This describes its image in the
divisorial log structure of $A$ and respects units and the structure
maps.

Pulling back the Kummer-\'etale cover gives a Kummer-\'etale cover of
this log trait. Proposition~A.10, now on $\operatorname{Spec}A$, says
that its generic extension is tame. It is the extension specified by
$V$ at $\eta$. Thus $V$ is tame at every codimension-one boundary point
of $\mathcal Y$, proving one inclusion of the displayed categories.
Interchanging the two models proves equality. Notice that the
valuative criterion supplied the required individual trait maps;
no proper regular common refinement was used.

For a model $\mathcal X$, let $j_{\mathcal X}$ and $i_{\mathcal X}$
be its interior and strict closed-fibre inclusions, and equip
\[
 E_{\mathcal X}=j_{\mathcal X}^*\circ(i_{\mathcal X}^*)^{-1}
\]
with the comparison isomorphisms of a chosen quasi-inverse. It is an
equivalence from the closed-log-fibre category to the common full
subcategory of $\operatorname{F\acute Et}(U)$. The comparison between
two models is $E_{\mathcal Y}^{-1}E_{\mathcal X}$, together with its
identification with the identity after restriction to $U$.
Full faithfulness gives a unique natural isomorphism between such
normalized comparisons; their composition is compatible for the same
reason. An actual morphism of models has the requisite commuting
restriction squares, so its pullback agrees with this comparison.
All these functors preserve pullback of covers of $S^{\log}$.

Applying the same source results to $S^{\log}$ identifies
\[
 \pi_1^{\mathrm{k\acute et}}(S^{\log})
 \simeq\pi_1^{\mathrm{k\acute et}}(s^{\log})
 \simeq G_F^t.
\]
The induced map $\Pi^t(U)\to G_F^t$ is surjective: a connected finite
tame extension of $F$ remains connected after pullback to the
geometrically connected $U$. This is the connected-cover criterion
for surjectivity of a morphism of fundamental groups.

Here is an explicit description of its framed section. Let $F^{\rm nr}$
be the maximal unramified extension, with residue field $k^{\rm sep}$.
The maximal tame extension is
\[
 F^t=\bigcup_{(m,p)=1}F^{\rm nr}(\varpi^{1/m}),
\]
with compatible roots; for $p=0$ there is no restriction on $m$.
Indeed, over the strictly henselian valuation ring in $F^{\rm nr}$,
a tame finite extension has ramification index $m$ invertible in the
residue field. Its uniformizer expresses $\varpi$ as a unit times an
$m$th power; Hensel's lemma extracts the root of the unit. This proves
the displayed description. All roots of unity of order prime to $p$
are in $F^{\rm nr}$ by the same lemma. Extend $G_k=\operatorname{Gal}
(F^{\rm nr}/F)$ by fixing the chosen roots of $\varpi$. This gives
$\sigma_{\varpi}$; conjugation on their root-of-unity multipliers is
the prime-to-$p$ cyclotomic action. This proves the semidirect-product
description and the stated exact sequence.

If $\varpi'=u\varpi$ with $u\in R^\times$, compatible roots of $u$
may be chosen in $F^{\rm nr}$ by lifting roots of $\overline u$ in
$k^{\rm sep}$. The new section differs by the inverse prime-to-$p$
Kummer cocycle of $\overline u$. Units congruent to $1$ have unique
such roots with residue $1$ and make no further ambiguity. Thus the
framing depends only on the first-order nonzero normal direction,
as in characteristic zero, but on the tame arithmetic quotient.
\end{proof}

\begin{proposition}[Refuted: preservation of the strict Isom-anabelian
 definition under reduction]\label{A6-frobenius-counterexample}
The characteristic-zero Isom-anabelian property used in the
compactification definition is not preserved unchanged under arithmetic
specialization. There are two nonisomorphic hyperbolic curves over
$k=\mathbf F_{16}$ whose augmented \'etale fundamental groups are
isomorphic over $G_k$. Both are good reductions of hyperbolic curves
over the number field $\Q(\zeta_{15})$.
\end{proposition}
\begin{proof}
Put
\[
 k=\mathbf F_2[a]/(a^4+a+1),\qquad
 U_a=\mathbf P^1_k\setminus\{0,1,\infty,a\},\qquad
 U_{a^2}=\mathbf P^1_k\setminus\{0,1,\infty,a^2\}.
\]
The quartic has no root in $\mathbf F_2$, and its remainder modulo
the only monic irreducible quadratic $T^2+T+1$ is $1$.
A reducible quartic has a factor of degree at most two, so it is
irreducible and $a$ has degree four over $\mathbf F_2$.

The map $z\mapsto z^2$ is a $k$-morphism $U_a\to U_{a^2}$:
on coordinate rings it fixes every element of $k$ and sends the
target coordinate $w$ to $z^2$. In particular
$w-a^2\mapsto(z-a)^2$. This is the relative Frobenius, not the
absolute Frobenius acting nontrivially on $k$. It is finite, radicial
and surjective, hence a universal homeomorphism. One can check these
properties directly from $z^2=w$: the affine map is finite of degree
two, it is integral and therefore surjective, and after every field
extension the equation has at most one root. Restricting away from
the displayed four-point sets preserves these properties.

The opened topological-invariance theorem
\cite[Proposition~58.8.4, Tag~0BQN]{A6StacksUniversalHomeomorphism}
therefore gives an equivalence of finite \'etale cover categories
and an isomorphism of their fundamental groups. Since the morphism
is over $k$, this isomorphism commutes with the augmentations to
$G_k$. This is a direct application of the source theorem, with its
universal-homeomorphism hypothesis verified above.

An isomorphism of the two open curves would extend uniquely to an
automorphism of their smooth proper completions $\mathbf P^1_k$,
and would permute the omitted four-point sets. Normalizing three
points shows that their fourth parameters would have to lie in the
same six-element cross-ratio orbit
\[
 \left\{a,\ 1-a,\ a^{-1},\ (1-a)^{-1},\
              a/(a-1),\ (a-1)/a\right\}.
\]
Equality of $a^2$ with any of these expressions would give a nonzero
polynomial over $\mathbf F_2$ of degree at most three annihilating
$a$, contrary to its degree four. Thus these are nonisomorphic
\emph{unlabelled} open curves; the obstruction is not an artifact of
adding puncture labels.

For the arithmetic lifts, $a$ has multiplicative order $15$:
$a^{15}=1$, while $a^3\ne1$ and $a^5\ne1$, as follows by reducing
powers using $a^4=a+1$. Thus evaluation $\zeta_{15}\mapsto a$
gives a prime $\mathfrak p$ over $2$ in $\Z[\zeta_{15}]$ with
residue field $k$. Indeed $\Phi_{15}(a)=0$, and the quotient of
$\Z[T]/(\Phi_{15}(T))$ by $(2,T^4+T+1)$ is precisely $k$.
The cyclotomic algebra is \'etale at this prime: its defining
polynomial divides $T^{15}-1$, whose reduction has no repeated root
in characteristic two. Hence its localization at $\mathfrak p$ is
an unramified discrete valuation ring, with fraction field
$K=\Q(\zeta_{15})$ and residue field $k$.

Over this ring the four sections $0,1,\infty,\zeta_{15}$ are
pairwise disjoint, and so are $0,1,\infty,\zeta_{15}^2$, because
their reductions are distinct. Their complements in $\mathbf P^1$
are smooth families with the displayed special fibres. The generic
fibres are smooth hyperbolic curves over the number field $K$, so
the smooth Isom theorem used in Theorem~\ref{A6-anabelian} applies to
them. The pair of special fibres violates that same strict Isom
property. Thus arithmetic specialization of cover categories must
not be promoted to preservation of that reconstruction property;
in positive characteristic the Frobenius ambiguity is already
present for genus zero and four punctures.
\end{proof}

\begin{proposition}[Refuted: replacing tame covers by all covers in
 mixed characteristic]\label{A6-wild-counterexample}
The first equivalence of Theorem~\ref{A6-tame-specialization} is false
with $\operatorname{F\acute Et}(U)$ in place of the tame subcategory,
even for a smooth proper relative curve with strict normal-crossings
boundary and hyperbolic generic interior.
\end{proposition}
\begin{proof}
Take $R=\Z_2$, $F=\Q_2$, $\mathcal X=\mathbf P^1_R$, and let its
boundary be the reduced union of the special fibre and the three
disjoint sections $0,1,\infty$. The total space is regular and proper.
At the intersections of the special fibre with a section the boundary
has equation $2z=0$ in a regular parameter system; elsewhere it has
one smooth component. Its divisorial log structure is therefore fs
and log regular, and its interior is
$U=\mathbf P^1_F\setminus\{0,1,\infty\}$.

Put $L=\Q_2(\alpha)$ with $\alpha^2=2$. The polynomial $T^2-2$ is
Eisenstein at $2$, so $L/F$ has degree two. It is separable in
characteristic zero. Consequently $V=U\times_F L\to U$ is a finite
\'etale cover; it is connected because $U$ is geometrically connected.
At the generic point of the special fibre, the valuation ring is
$A=\Z_2[z]_{(2)}$, with residue field $\mathbf F_2(z)$.
The same Eisenstein polynomial gives the extension $A[\alpha]$.
This is a noetherian local domain with maximal ideal $(\alpha)$:
$2=\alpha^2$ and its quotient is $\mathbf F_2(z)$. It is a
one-dimensional regular local ring, hence a discrete valuation ring
and the normalization of $A$ in $F(z)(\alpha)$. Its ramification
index is exactly two, since $v(\alpha)=1$ and $v(2)=2$.
This is wild in residue characteristic two. Hence $V$ is absent from
$\operatorname{Tame}_{\mathcal X}(U)$ and cannot extend to a finite
Kummer-\'etale cover of $\mathcal X^{\log}$, by Proposition~A.10.
The same constant extension already gives the obstruction on the
base log trait. In particular an identification with the full $G_F$
would discard a real obstruction, not merely a choice of notation.
\end{proof}

\paragraph{The remaining model-existence scope.}
The theorem compares every \emph{existing} proper log-regular model,
in arbitrary dimension, and it compares any two such models without
requiring a common log-regular refinement. It does not assert existence
of such a model for an arbitrary mixed-characteristic generic variety,
nor cofilteredness of their geometric model category. The
characteristic-zero resolution theorem used in
Theorem~\ref{A6-model-specialization} has characteristic-zero
\emph{total-space} hypotheses and does not prove either assertion over
$\Z_p$. Altering the generic open by a generically finite cover would
also change the problem. No such mixed-characteristic resolution or
chosen-special-fibre reconstruction assertion is used here.

The parenthetical cyclotomic assertion on 67/049(b) has a separate
qualification: a rational node need not have rational branches.

\begin{proposition}[Refuted without split branches; exact corrected
 inertia character]\label{A6-node-character}
Let $e$ be a geometric node of a stable curve in the framed logarithmic
interpretation of Theorem~\ref{A6-anabelian}, over a characteristic-zero
field $K$. On the subgroup of $G_K$ stabilizing $e$, let
$\varepsilon_e$ be the character which is $1$ when the two branches
are preserved and $-1$ when they are interchanged. The character on
its node inertia group is
\[
 I_e\simeq\widehat{\Z}(1)\otimes\varepsilon_e,
 \qquad \gamma\longmapsto\chi_{\rm cyc}(\gamma)\varepsilon_e(\gamma).
\]
In particular the invariant-conjugacy-class condition does not by
itself make the character cyclotomic. It becomes cyclotomic after
passing to the branch-splitting field. Marked-point inertia has the
ordinary cyclotomic character on its point stabilizer.
\end{proposition}
\begin{proof}
After splitting the branches, choose the local smoothing chart
$xy=c\tau$ with its fixed normal frame. Kummer inertia on the
$x$-branch acts on an $m$th root of $x$ by a root of unity $\zeta_m$.
The admissible identification with the $y$-branch is inverse:
on compatible roots it acts by $\zeta_m^{-1}$. An arithmetic element
preserving the branches therefore acts on inertia by
$\chi_{\rm cyc}$, by its action on roots of unity. An element
interchanging the branches introduces the additional inverse,
hence the displayed factor $\varepsilon_e$. Passing to inverse limits
proves the formula. It does not depend on the branch coordinates or
on an orientation chosen to identify $I_e$ with $\widehat\Z$.
The inertia subgroup is self-normalizing in the geometric group, by
\cite[Proposition~1.2(ii), p.~461]{Mochizuki2007A567}; thus choosing a
representative of its invariant conjugacy class introduces no extra
character. At a marked point there is only one branch, so no sign
occurs.

For a completely explicit counterexample over $\Q$, take the projective
plane cubic
\[
 Y^2Z=X^2(X-Z),
\]
marked at $[0:1:0]$. Its unique singularity is $[0:0:1]$, an ordinary
node with tangent cone $y^2+x^2=0$. To check uniqueness, the affine
critical equations are $2y=0$ and $x(2-3x)=0$; the candidate
$(2/3,0)$ does not lie on the cubic, since
$y^2-x^3+x^2=4/27$ there. At infinity the partial derivative with
respect to $Z$ is $1$. The normalization is $\mathbf P^1$ with
parameter $a=y/x$, and map
\[
 x=a^2+1,\qquad y=a(a^2+1).
\]
Homogeneously it is
$[A:B]\mapsto[B(A^2+B^2):A(A^2+B^2):B^3]$; it has no base point
and is birational. The two node branches are $a=i,-i$, and the
marking is $a=\infty$. Thus the one-pointed curve is geometrically
stable of type $(1,1)$.

The normal tensor
\[
 N=\partial_a|_{i}\otimes\partial_a|_{-i}
\]
is nonzero and Galois invariant: the nontrivial element of
$\operatorname{Gal}(\Q(i)/\Q)$ interchanges its two factors, and the
normal tensor uses their symmetric identification. The node itself is
rational. Complex conjugation has $\varepsilon_e=-1$ and
$\chi_{\rm cyc}=-1$, so it acts as $+1$ on $I_e$, whereas the
cyclotomic character would act as $-1$. This verifies every
hypothesis and refutes the automatic-cyclotomic parenthesis on the
sheet. It does not alter the reconstruction theorem: node \emph{Dehn
twists} and branch inertia are different groups, and the twist action
used in Lemma~\ref{A6-normal-kummer} is the already cited intrinsic
cyclotomic twist action.
\end{proof}
